\section{The circular collision hierarchy}
\label{sec:results}
A time window which admits two collisions need not test dependence between
successive complete free flights. Under stationary preparation, its count
law can depend only on the marginal law of one flight. The onset of higher
collision tails is different: it requires an actual orbit segment, and its
probability measures how a single two-dimensional collision state shadows
a prescribed sequence of reflections. We determine this hierarchy for a
periodic Lorentz gas, including its dependence on the length of the segment.

Let $\Lambda=\Z(1,0)+\Z(1/2,\sqrt3/2)$ and remove the closed radius-$R$
disk from $\R^2/\Lambda$. The particle has unit speed and reflects
specularly. Throughout,
\begin{equation}\label{eq:parameters}
 0<R_-\le R\le R_+<\tfrac12,\quad K=[R_-,R_+],\quad
 g_R=1-2R,\quad A_R=\tfrac{\sqrt3}{2}-\pi R^2,\quad
 \lambda_R=\frac{2R}{A_R}.
\end{equation}
Preparation is normalized phase volume $\dd q\dd\vartheta/(2\pi A_R)$.
Write $N_t$ for the number of collisions in $(0,t]$. The outgoing collision
section has angular coordinates $x=(\alpha,\phi)$, with
$|\phi|<\pi/2$, and invariant probability
\begin{equation}\label{eq:nu}
 \dd\nu=\frac{\cos\phi}{4\pi}\dd\alpha\dd\phi.
\end{equation}
Its collision map is $B_R$, its physical roof is $\tau_R$, and
$S_j\tau_R=\sum_{i=0}^{j-1}\tau_R\circ B_R^i$. These are the actual
mechanical map and roof. In particular, the summands are not independent.
Finite horizon is not required: the proof uses the positive lower bound
$\tau_R\ge g_R$, not an upper bound.

Put
\begin{equation}\label{eq:chi}
 c_R=\frac{1-R}{R},\qquad
 \chi_R=\arcosh c_R,\qquad s_R=\sinh\chi_R=\frac{\sqrt{g_R}}{R},
 \qquad \chi_* =\min_{R\in K}\chi_R>0.
\end{equation}
There are six unit lattice vectors $e$. Let $x_e$ denote the outgoing
normal state facing $e$; then $B_Rx_e=x_{-e}$. In angular coordinates
these states do not depend on $R$.

\subsection{Unweighted threshold law}
\begin{theorem}[Uniform onset at every collision order]\label{thm:unweighted}
There are $\varepsilon_0>0$ and $C<\infty$, depending on $K$ but not on
$j$, such that $\varepsilon_0<\min_K g_R/4$ and, for all $j\ge1$,
$R\in K$ and $0<\varepsilon<\varepsilon_0$,
\begin{equation}\label{eq:main-onset}
 \Pr_R\{N_{jg_R+\varepsilon}\ge j+1\}
 =\Pr_R\{N_{jg_R+\varepsilon}=j+1\}
 =C_j(R)\varepsilon^2\bigl(1+\varepsilon h_j(R,\varepsilon)\bigr),
 \qquad C_j(R)=\frac{3\lambda_R}{2R\sinh(j\chi_R)}.
\end{equation}
The functions $h_j$ are smooth up to $\varepsilon=0$. Every fixed
mixed derivative of $h_j$ in $R$ and $\varepsilon$ is bounded uniformly
in $j$ and $R\in K$, after decreasing $\varepsilon_0$ if necessary.
For $-\varepsilon_0<\varepsilon\le0$ the probability in
\eqref{eq:main-onset} is zero.
\end{theorem}

The theorem concerns the entire equilibrium event, not a prescribed
transverse initial subset. Its threshold collar has width independent
of the collision order. The observation times $jg_R+\varepsilon$ grow
with $j$; the events remain the maximal possible collision counts in
those windows. This is a uniform extremal regime, not a statement about
typical long-time fluctuations.

\subsection{Sources and summable response}
Let $f_R$ be a smooth real-valued function of the outgoing collision
state. It is enough that it be smooth on fixed neighborhoods of the six
normal states, uniformly on a neighborhood of $K$; choose any bounded
extension elsewhere. Let $F_*\ge1$ bound $|f_R|$ on those neighborhoods.
On the event in \eqref{eq:main-onset}, write $X_0,\ldots,X_j$ for the
outgoing states at its successive impacts and set
\[
 Z_j(R,q,\varepsilon)=
 \E_R\left[\one_{\{N_{jg_R+\varepsilon}=j+1\}}
             \exp\left(q\sum_{i=0}^j f_R(X_i)\right)\right],
 \qquad
 w_{j,e}(R)=\sum_{i=0}^j f_R(x_{(-1)^i e}).
\]
The source is attached to a correlated impact record, not independently
to a list of sampled roofs.

\begin{theorem}[Marked threshold expansion]\label{thm:marked}
For every fixed $q_0<\infty$, there are smooth functions
$h_{j,e}(R,q,\varepsilon)$, uniformly bounded with every fixed mixed
derivative for $R\in K$, $|q|\le q_0$ and
$0\le\varepsilon<\varepsilon_0$, such that
\begin{equation}\label{eq:marked-onset}
 Z_j(R,q,\varepsilon)=
 \frac{\lambda_R\varepsilon^2}{4R\sinh(j\chi_R)}
 \sum_{|e|=1} e^{q w_{j,e}(R)}
           \bigl(1+\varepsilon h_{j,e}(R,q,\varepsilon)\bigr).
\end{equation}
All constants are independent of $j$. The offset $\varepsilon$ is held
fixed in radius derivatives in this formula. If $q_0F_*<\chi_*$, then
for each fixed $a,b,c\ge0$,
\begin{equation}\label{eq:summable}
 \sup_{\substack{R\in K,\ |q|\le q_0\\0\le\varepsilon<\varepsilon_0}}
 \left|\partial_R^a\partial_q^b\partial_\varepsilon^c
       \frac{Z_j(R,q,\varepsilon)}{\varepsilon^2}\right|
 \le C_{a,b,c}(1+j)^{a+b}
           e^{-(\chi_*-q_0F_*)j}.
\end{equation}
The quotient is understood by its smooth right extension at zero.
Consequently its series over $j$ and every fixed mixed derivative
converge uniformly on the displayed set. The same conclusions hold
for finitely many sources, with $qf$ replaced by a scalar product and
$q_0F_*$ by a uniform bound for its absolute value per impact.
\end{theorem}

Uniformity in $j$ is the main issue. An initial-value differentiation
of $B_R^j$ would amplify errors exponentially. Instead we parametrize
an orbit segment by the transverse positions of its first and last
impacts. Its effective action remains uniformly nondegenerate, while
the flux density has size $\sinh(j\chi_R)^{-1}$. A relative determinant
estimate, rather than an absolute error estimate, preserves this scale.

\subsection{Consequences and relation to the fixed-window theorem}
Theorem~\ref{thm:marked} gives the finite-order response at all maximal-count
onsets, including the derivative jumps across them. Section~\ref{sec:records}
retains collision times, positions, both velocities and sampling-time cuts.
Section~\ref{sec:consequences} compares the hierarchy with independent roofs
having exactly the same marginal and recovers the normal periodic-orbit
multiplier from the coefficients $C_j$.

The accompanying article \cite{TC} retains the full proof of the
$T=11/100$, $R\in[9/20,47/100]$ three-point count law and its regular-level
response. That is a complementary window: its second-collision probability
is already positive throughout the radius interval. Here the count changes
support at $t=jg_R$, and its second response is singular there. The
one-flight onset coefficient suggested by that calculation is the case
$j=1$ of \eqref{eq:main-onset}; the uniform multi-flight determinant and
the full marked record follow from the uniform boundary-value argument
and the explicit specialization below.
